\section{Stakeholders and Operational Use Cases}
The client application supplies a prompt and receives either a previously completed answer, a newly generated answer, or an explicit error. The developer integrating that client needs a stable single-turn API, visible outcome metadata, and a way to request private or no-store treatment. An organization administrator needs control over credential membership, workspace permission, provider profiles, and retention limits. An operator needs readiness, counters, and failure reasons to tell an empty cache from an unhealthy verifier or provider. The person whose content appears in a request needs the gateway to keep private material outside shared workspace entries. These interests constrain the same request path; optimizing hit rate alone would not satisfy them.

The primary use case begins when an authenticated application submits a supported request. The gateway resolves the principal and provider route, classifies the requested scope, and checks whether the text can be cached. If an exact answer exists under that identity, it may return it after a current-entry check. Otherwise semantic mode may retrieve a bounded shortlist and assess each stored answer against the new task. An accepted and still-current candidate gives a HIT. If the shortlist is empty, rejected, uncertain, conflicting, or over budget, the gateway calls the configured provider. A normally completed JSON response is considered for storage only if retention policy permits, then returned with its storage outcome. Provider or admission failure produces an error rather than an invented cached answer.

A second use case is authorized collaboration. Alice asks a reusable workspace question, and Carol later asks an equivalent question as another member of that same workspace. The answer may be reused only when Alice's content was eligible for shared storage, Carol has permission, and the provider compatibility identity agrees. In the selected isolated HTTP probe, a seeded answer was served to Carol; another valid paraphrase of Alice's completed response missed despite being in scope. This distinction makes workspace authorization a precondition, not a promise that all semantically related requests will hit.

A third use case is privacy or uncertainty. An explicit private request must remain in a subject/application scope. A secret-bearing or volatile request can bypass storage. A changed technical action or target should decline an old answer and let the provider handle the new task. If a verifier cannot decide within its token or time bound, the safe outcome is a MISS. This can increase cost and latency; the report measures such lost opportunities rather than hiding them inside an overall success count.

\section{Inputs, Outputs, and Acceptance Criteria}
The application supplies a bounded JSON body, a gateway credential, and optionally a scope request. A prompt or one supported user message is required. A workspace identifier is meaningful only after the resolved principal proves membership and permission. Extra fields are rejected by the current schema rather than silently forwarded. A successful JSON response identifies the outcome and includes generated or reused content; a streaming response has content events followed by completion metadata only after normal termination. Metrics and candidate traces are operational evidence, not a substitute for client-visible correctness.

Acceptance of a cached answer has three separate conditions. First, the principal and requested scope must authorize access to that entry. Second, the candidate must still be present and unexpired when selected. Third, a semantic candidate must pass the selected suitability policy; high vector similarity alone is insufficient. The exact shortcut uses request identity and entry checks rather than running NLI. A candidate may pass retrieval and fail suitability, as in the saved live challenges. Conversely, a policy might approve an unsuitable candidate on an authored test. The test plan therefore records retrieval inclusion, policy action, request HIT/MISS, and independent label separately.

For a newly generated answer, successful completion requires a configured provider route to return a normal terminal signal within its limits. The gateway can deliver a completed answer even if a later cache write fails, while reporting that failure. It must not cache an incomplete stream or an upstream error. For an authorized workspace request, the response may be narrowed to private storage if the content classification requires it. These are observable behaviors that can be asserted with local mock providers; they do not establish production answer quality for a real external provider.

\section{Nonfunctional Constraints and Threat Boundaries}
Latency has multiple budgets. Authentication and admission precede the gateway timer reported in a JSON response. Cache lookup can include Redis reads, a local embedding, ANN search, and NLI inference. A miss additionally pays the configured provider's latency. The selector limits candidate count and verification time to avoid letting an unusually dense shortlist consume an unbounded critical path. A finite limit can convert a potentially reusable answer into a MISS; this is a documented trade-off rather than an invisible timeout.

Scope separation is enforced through server-resolved identity and an opaque Redis tag. An attacker with an application credential must not gain another workspace by submitting a workspace field. A same-worded question in another organization must use a different cache namespace. The content classifier is a defense for obvious private markers, but arbitrary confidential text may evade heuristics. The deployment therefore still depends on correct client classification, credential governance, Redis protection, and appropriate retention settings. The local tests cover selected paths, not a complete adversarial privacy assessment.

Capacity limits include entry count, bytes, TTL, and least-recently-used eviction at scope and organization levels. These protect finite storage but can lower hit rate under pressure. Multiple gateway workers share Redis state, while each holds its own models and local coalescing state. Memory scales with worker count; duplicate provider calls across workers remain possible. A daily admitted-miss counter is an operational guard and cannot be interpreted as a billed-cost ceiling because token lengths and provider prices vary.

\section{Traceability and Unresolved Release Criteria}
The source map in Chapter~5 and the case catalogue in Chapter~6 connect requirements to implementation and evidence. Authentication and scope have test assertions and selected HTTP observations. Semantic suitability has labelled development and frozen challenges, with known valid misses and some unsafe approvals. Provider completion and storage failure have injected-failure tests. The dashboard's final browser behavior, exact-final-revision full suite, independently reviewed answer labels, real-provider billing, and target-scale operation remain release criteria. A historical pass count cannot close a criterion after runtime code changes, and a controlled provider cannot establish external provider economics.

\section{Numbered Requirements and Feasibility}
The requirements below use ``shall'' for implemented contract behavior at the recorded checkpoint. They do not state a numeric production service-level objective. The table in the preceding section maps major requirements to components and saved checks.

\begin{enumerate}
\item[FR1.] The gateway shall validate a single-turn request and authenticate a gateway credential before accessing scoped cache content.
\item[FR2.] It shall resolve organization, subject, application, workspace permission, and provider profile on the server; an override shall not grant membership.
\item[FR3.] It shall search only entries under a compatible scope and configured policy, with an exact shortcut and an optional bounded semantic shortlist.
\item[FR4.] When answer-aware mode is enabled, it shall assess the stored answer for the incoming task and abstain on unsupported, conflicting, or over-budget candidates.
\item[FR5.] It shall call an explicitly configured provider on a miss or bypass and require a normal complete response before considering storage.
\item[FR6.] It shall honor no-store and private restrictions, TTL and quota limits, and shall not turn a cache-write failure into loss of a completed provider answer.
\item[NFR1.] Lookup work shall have finite candidate and time bounds, and request/provider deadlines shall yield explicit outcomes.
\item[NFR2.] Shared state shall be scoped and observable through outcome metadata, timing and reason traces without exposing provider secrets in cache keys.
\item[NFR3.] Readiness shall depend on configured required components; multiple workers shall not assume process-local coalescing is shared.
\end{enumerate}

Technical feasibility is supported by the local Compose deployment, the saved Redis Stack and model-backed HTTP probe, and bounded model resource measurements. It does not establish the capacity of a production tenant population. The pinned NLI model's observed resident-memory increase of about 681 MB in one process makes worker count a deployment constraint. Redis Stack must be provisioned with a compatible vector index and finite memory; a semantic profile also needs local model artifacts. The mock profile permits repeatable integration without paid calls, while an external provider profile needs a network route and server-held credentials. Economic feasibility remains conditional: expected provider generation avoided by suitable HITs must exceed gateway inference, Redis, hosting, and operational cost. The stress study counts controlled-provider calls but provides no billed token or infrastructure totals, so the break-even inequality in Chapter~6 is a decision model rather than a savings result.
